\section{Free critical cooling preserves a small target component}
\label{sec:critical-free-cooling}

Write $F(q)=\mathbb E\tanh^2(\sqrt qZ)$ for a standard Gaussian $Z$.
For a feature second moment $Q$, let $\mathsf M_Q$ denote the law of
coordinatewise tanh of a centered Gaussian with covariance $Q$.

\begin{lemma}[Uniform free-cooling estimates]
\label{lem:critical-free-cooling}
There are universal constants $B<\infty$ and $b>0$ with the following
property. Suppose $Q_0$ has constant diagonal $q_0\in[F(1),F(2)]$ and
$Q_0\succeq v_0ff^T$, where $f$ is Boolean and $v_0>0$.
Successively apply the unmodified feature law $\mathsf M_Q$.
For every $0<\delta\le1/128$, the first layer whose diagonal $q_k$
is at most $\delta$ occurs after at most $B/\delta$ transitions.
At this layer,
\[
 \delta/2\le q_k\le\delta,
 \qquad Q_k\succeq b v_0\delta\,ff^T.
\]
All these transitions cost zero in the critical variational functional.
\end{lemma}

\begin{proof}
The diagonal evolves exactly by $q_{j+1}=F(q_j)$.
The function $F$ is increasing and $0<F(q)<q$ for $q>0$.
Taylor expansion with bounded derivatives gives
\[
 F(q)=q-2q^2+O(q^3).
\]
It follows by continuity on $(0,F(2)]$ and the limit at zero that
there is a universal $c_F>0$ satisfying
$q-F(q)\ge c_Fq^2$ throughout this interval.
Hence $1/q_{j+1}-1/q_j\ge c_F$, which bounds the first hitting time
of $(0,\delta]$ by $O(1/\delta)$.
The elementary bound $F(q)\ge q/(1+3q)$ follows from
$\tanh^2t\ge t^2/(1+t^2)$ and Cauchy--Schwarz.
Monotonicity therefore implies that the first value below $\delta$
is at least $F(\delta)\ge\delta/(1+3\delta)\ge\delta/2$.

If $Q\succeq vff^T$ has constant diagonal $q$, decompose its Gaussian
as $\sqrt v\,fZ+G$, with $G$ independent and covariance $Q-vff^T$.
The conditional feature mean is $f\psi(Z)$, with the same scalar
function at every input because all residual variances equal $q-v$.
Conditional expectation, projection onto the first Hermite polynomial,
and Gaussian integration by parts give
\[
 Q'\succeq\mathbb E\psi(Z)^2ff^T
       \succeq v[1-F(q)]^2ff^T.
\]
Thus we may track $v_{j+1}=v_j[1-F(q_j)]^2$.
The ratio $v_j/q_j$ is multiplied by
\[
 R(q_j)=\frac{q_j[1-F(q_j)]^2}{F(q_j)}.
\]
The expansion of $F$ yields $R(q)=1+O(q^2)$ at zero, and $R$ is
positive on $(0,F(2)]$. Consequently
$|\log R(q)|\le C_Rq^2$ there for a universal $C_R$.
Also
\[
 \sum_{j<k}q_j^2
 \le c_F^{-1}\sum_{j<k}(q_j-q_{j+1})\le F(2)/c_F.
\]
The product of all ratio factors is therefore at least a universal
positive constant. Since $q_0<1$ and $q_k\ge\delta/2$, this gives
$v_k\ge b v_0\delta$ for a universal $b>0$.
\end{proof}
